% Study sheet: Volume I, Chapter 16, Conclusion.
% Numbers: \val{...} from conclusion_studysheet_numbers.tex, written by conclusion_studysheet_numbers.py.
\documentclass{studysheet}
\input{conclusion_studysheet_numbers}
\sheetvolume{1}
\sheetchapter{16}{Conclusion}
\sheetsources{Numbers from the mlsysim registry and the chapter's scenarios; floors are idealized lower bounds; assumed inputs are marked.}

\begin{document}
% ===================================================================== FRONT: LEARN
\frontside
\begin{bigidea}
\hl{The system is the model:} pipeline, weights, execution stack, serving path and monitoring together
produce what a user receives. Two threads run through the volume: the \hl{D·A·M taxonomy}, since decisions
couple across Data, Algorithm and Machine and the costliest failures happen at their interfaces, and the
\hl{iron law}, which says where the time goes. What lasts is the method: name the limiting term, measure it, budget for it.
\end{bigidea}

\begin{leftcol}
\sheetsec{Key ideas}
\begin{ideas}
\idea{Test the interfaces.}{Compare two paths' outputs on the same inputs (skew); track each cohort.
  Roll back a bad release; fix a shift where it entered.}
\idea{A regime belongs to a workload on a machine.}{Batch, phase, precision and machine move it:
  MobileNetV2 falls below an H100's higher ridge.}
\idea{Aim the lever at the limiting term.}{Shrinking a term that does not limit buys little.}
\idea{Principles differ in kind.}{Bounds, decompositions, fitted models, policies and heuristics;
  each holds only within its stated assumptions.}
\idea{Complexity is displaced, not removed.}{After a local simplification, trace where the cost
  landed: quantizing a model moves work into validation.}
\idea{The dual mandate is a budget kept.}{Trust (drift, skew, bias feedback) and delivery (time,
  memory, energy) are measured continuously.}
\end{ideas}

\vspace{10pt}
\sheetsec{The fleet case: one failure, every team passed}
{\centering
\begin{tikzpicture}[sheet, every node/.style={nbox, text width=44pt, align=center, minimum height=18pt}]
  \node (e) at (26,16) {evaluation path};
  \node (c) at (88,16) {INT8 calibration};
  \node (f) at (26,-16) {firmware path};
  \node (s) at (88,-16) {cohort inputs shift};
  \node[accbox, text width=40pt, align=center, minimum height=18pt] (m) at (153,0) {quantized model};
  \node[hotbox, text width=44pt, align=center, minimum height=18pt] (d) at (220,0) {cohort accuracy drops};
  \draw[arr] (e) -- (c); \draw[arr] (f) -- (s); \draw[arr] (c) -- (m); \draw[arr] (s) -- (m); \draw[arr] (m) -- (d);
\end{tikzpicture}\par}

\vspace{12pt}
\sheetsec{Lighthouses: what limits, what pays}
\begin{sheettable}{@{}l l L L@{}}
\thead{Lighthouse} & \thead{Regime} & \thead{Lever that pays} & \thead{Wasted}\\ \midrule
ResNet-50 & compute & batching, low precision & bandwidth\\
LLM decode & bandwidth & batching, fewer bytes/weight & peak math\\
DLRM & capacity & shard and cache tables & dense math\\
MobileNetV2 (phone) & compute & depthwise, then INT8 & op cuts once bandwidth-bound\\
Keyword spotting & power & tiny INT8 model & cloud streaming\\
\end{sheettable}
\end{leftcol}%
\begin{rightcol}
\sheetsec{Three questions, thirteen principles}
\begin{sheettable}{@{}l L@{}}
\thead{Question} & \thead{Principles}\\ \midrule
What shapes predictions & data as code; bias feedback\\
What limits execution & data gravity; iron law; silicon contract; Pareto frontier; arithmetic intensity; energy movement; Amdahl\\
What supports deployment & verification gap; drift; training-serving skew; latency budget\\
\end{sheettable}

\vspace{12pt}
\sheetsec{Formulas you will use}
\fcard{Iron law, and which floor wins}%
  {\Time = \under{\DataVolume/\Bandwidth}{data} + \under{\Operations/(\PeakRate\,\EfficiencyHw)}{compute}
   + \under{\Latency}{overhead}}%
  {Bytes $\DataVolume$ at bandwidth $\Bandwidth$, operations $\Operations$ at peak $\PeakRate$ and efficiency
   $\EfficiencyHw$, fixed overhead $\Latency$; without overlap the terms add. Ideal floors ($\EfficiencyHw = 1$):
   $\IronLawMemoryTime/\IronLawComputeTime =$ ridge $\PeakRate/\Bandwidth$ $\div$ intensity $\Operations/\DataVolume$.
   Batch-one 16-bit decode: \val{fpp} FLOPs per \val{bpp}-byte weight, intensity \val{llint}.}%
  {D·A·M form: samples per task (Data) $\times$ FLOPs per sample (Algorithm) $\times$ seconds per FLOP (Machine).}{}
\fcard{Amdahl's law}%
  {\text{Speedup} = \frac{1}{(1-\ParallelFraction) + \ParallelFraction/\AcceleratedFractionSpeedup}}%
  {Share $\ParallelFraction$ sped up $\AcceleratedFractionSpeedup$ times; the rest caps the gain.}{}{}
\fcard{Latency budget}%
  {\Pr[\EndToEndRequestLatency \le \CompletionDeadline] \ge \TargetCompletionQuantile}%
  {Critical path: serial stages add, parallel ones count only the slowest. The product sets
   $\CompletionDeadline$ and $\TargetCompletionQuantile$: a tail, not a mean.}{}{}
\fcard{From node to fleet}%
  {\MeanTimeBetweenFailures \approx \ComponentMttf / \AcceleratorCount}%
  {$\AcceleratorCount$ independent accelerators fail $\AcceleratorCount$ times as often; a synchronous job
   halts at any one, so it checkpoints.}{}{}
\end{rightcol}

% ===================================================================== BACK: PRACTICE
\backside
\begin{answers}
\ans{Finish}{\val{gxgb} GB $\div$ \val{gxbw} TB/s $=$ \val{gxtmem} ms; \val{gxgflop} GFLOP $\div$ \val{gxpeak} TFLOP/s
  $=$ \val{gxtcmp} ms; ratio $=$ ridge \val{gxpeak} $\div$ \val{gxbw} $\approx$ \val{gxridge} over intensity \val{llint}, as before:
  both floors shrank together. Cut bytes per weight.}
\ans{1}{(a) Capacity: \val{dlrmgb} GB $>$ \val{gpugb} GB; shard and cache tables; faster math is
  wasted. (b) Below the H100's higher ridge its
  depthwise layers turn bandwidth-bound: cut bytes, not operations. (c) Data (half the bytes per weight); overhead $\Latency$.}
\ans{2}{\val{q2retr} + \val{q2plan} + \val{q2ta} + \val{q2gen} + \val{q2ver} + \val{q2orch} $=$ \val{q2path} ms (not \val{q2seq}:
  the tools overlap), \val{q2slack} ms under the deadline. The \val{q2tb} ms tool is off the critical path, so speeding it up buys nothing.}
\ans{3}{The skew check: the same captured images through both paths, outputs compared; accuracy per cohort. No rollback: an older model sees the same shifted inputs; fix the path where the shift entered.}
\ans{4}{\val{mttfh} h $\div$ \val{q4n} $=$ \val{q4h} h between failures; \val{q4runh} h $\div$ \val{q4h} h $\approx$ \val{q4fail}
  stops: checkpoint and restart.}
\end{answers}

\begin{leftcol}
\sheetsec{Worked example}
\begin{practicebox}\boxtitle{What limits one token of Llama 2 \val{llparams}B decode?}
Batch one, 16-bit weights, sharded across \val{llshards} H100s: \val{llbw} TB/s and \val{llpeak} TFLOP/s in total.
Two FLOPs per parameter.\par\vspace{5pt}
{\centering
\begin{tikzpicture}[sheet]
  \draw[dim] (10,0) -- (222,0);
  \foreach \tk/\tl in {0.1/0.1,1/1,10/10,100/100,1000/1000}{
    \pgfmathsetmacro\x{10+(log10(\tk)+1)*53}
    \draw[dim] (\x,0) -- (\x,-2.5);
    \node[small, anchor=north] at (\x,-4) {\tl};}
  \node[small, anchor=north east] at (222,-13) {FLOP per byte};
  \pgfmathsetmacro\xi{10+(log10(\val{llint})+1)*53}
  \pgfmathsetmacro\xb{10+(log10(\val{ridge})+1)*53}
  \fill[accentsoft] (\xi,-1.5) rectangle (\xb,1.5);
  \draw[accent, dashed, line width=0.6pt] (\xb,-3) -- (\xb,16);
  \node[note, anchor=south] at (\xb,17) {ridge \val{ridge}};
  \filldraw[fill=accent, draw=accent] (\xi,0) circle (2.4pt);
  \node[note, anchor=south] at (\xi,5) {decode: \val{llint}};
  \node[small, anchor=south] at ({(\xi+\xb)/2},3) {bandwidth-bound side};
\end{tikzpicture}\par}
\begin{steps}
\item Bytes and work: \val{llparams}B parameters $\times$ 2 bytes $=$ \val{llgb} GB read; \val{llgflop} GFLOP.
\item Memory floor: \val{llgb} GB $\div$ \val{llbw} TB/s $=$ \res{\val{lltmem} ms}. Compute floor:
  \val{llgflop} GFLOP $\div$ \val{llpeak} TFLOP/s $=$ \res{\val{lltcmp} ms}.
\item Ratio \res{\val{llratio}$\times$} $=$ ridge \val{llpeak} $\div$ \val{llbw} $\approx$ \val{ridge} over intensity \val{llint} FLOP/byte.
\end{steps}
\takeaway{Compute tuning cannot lower the \val{lltmem} ms floor. Fewer bytes per weight can; batching shares
each read across requests, raising throughput.}
\end{practicebox}

\vspace{12pt}
\sheetsec{Your turn: finish it}
\begin{practicebox}\boxtitle{Does a model \val{scale}$\times$ smaller change the regime?}
Same method, now GPT-2 XL (\val{gxparams}B parameters) on one H100, batch one, 16-bit: \val{gxbw} TB/s, \val{gxpeak} TFLOP/s.\par\vspace{4pt}
\textbf{Given:} each token reads \val{gxgb} GB and does \val{gxgflop} GFLOP.\par
\textbf{You:} find both floors and their ratio, explain why the ratio did not move, and name the lever that pays.
\end{practicebox}
\end{leftcol}%
\begin{rightcol}
\sheetsec{Common mistakes}
\mistake{``We made our stage \val{amlocal}$\times$ faster, so the request is much faster.''}%
  {Inference was \val{amfrac}\% of the request. Amdahl gives
   \val{amsys}$\times$. Profile end to end.}{}
\mistake{``One accuracy number tells us the model is good.''}%
  {The fleet passed its aggregate test while one cohort lost \val{fleetdrop} points. Under a \val{slo} ms
   tail SLO, \val{accA}\% at \val{latA} ms cannot ship; \val{accB}\% at \val{latB} ms meets it.}{}
\mistake{``Faster arithmetic is what saves energy.''}%
  {One DRAM access (\val{edram} pJ) costs \val{elow}$\times$ an FP32 operation (\val{efpthirtytwo} pJ) and \val{ehigh}$\times$
   an FP16 one (\val{efpsixteen} pJ). Count both: where accesses outweigh the arithmetic, cut traffic.}{}

\vspace{2pt}
\sheetsec{Check yourself}
\begin{qs}
\item Name the regime and the lever that pays: (a) a DLRM ranker with \val{dlrmgb} GB of tables, \val{gpugib} GiB
  (\val{gpugb} GB) per accelerator; the team wants faster math.
  (b) MobileNetV2 moved from a phone NPU to an H100. (c) In batch-one decode, which iron-law term do INT8
  weights (from 16-bit) shrink, and which does removing Python dispatch shrink?
\item \assumed{} a composed request: retrieval \val{q2retr} ms, planner \val{q2plan} ms, two tools in parallel
  (\val{q2ta} and \val{q2tb} ms), generation \val{q2gen} ms, verifier \val{q2ver} ms, orchestration \val{q2orch} ms;
  deadline \val{q2dead} ms. Does it fit? Which speedup buys nothing?
\item In the fleet case, which measurement would have caught the drop before users did? Should the
  team roll back the model?
\item \assumed{} \val{q4n} accelerators, each with a \val{mttfh}-hour MTTF, run a synchronous \val{q4days}-day job.
  How often does the pool fail, how many stops hit the run, and what must the job do?
\end{qs}
\end{rightcol}
\end{document}
